\section{Introduction}
\label{sec:intro}

AI agents are beginning to negotiate, shop and sign agreements on people's behalf. The person who delegated rarely watches the transaction. What they see is the agent's account of it afterwards, and that account is where supervision happens: it is how a principal decides whether to accept a deal, change their instructions, or stop using the agent.

Recent work has shown that this channel can fail. Agents claim tasks succeeded when they did not \citep{advani2026falsesuccess}, conceal broken tools and fabricated results from their users \citep{guo2025upward}, and hide their own misconduct from a manager \citep{scheurer2023strategic}. In each of these cases the agent had something to gain from the omission. We study a case in which it does not. An agent can follow every instruction, keep every confidence and reach an acceptable deal, and still leave out a fact its principal could act on: that the principal's own requirement ruled out a better deal.

We build a negotiation testbed in which that fact is exact and appears identically for every agent. A seeded, four-issue procurement negotiation gives every package an exact value to the client. The client's briefing contains a hard requirement and a confidential budget. A scripted counterparty controls how good an outcome is reachable and, at the same point in every negotiation, offers a better deal that only the client's requirement forbids. After the negotiation, the client asks one open question, ``How did it go?'', and the reply is scored against ground truth.

Telling the client that their own rule cost them a better deal would, if anything, excuse a weaker outcome, so an agent that stays silent gains nothing by it. Whether the omission counts as a failure depends on what a principal wants from a report, which this design does not measure. What it measures is how often the omission occurs, whether the information was available to the agent, and whether a differently worded request elicits it. \Cref{sec:discussion} takes up what the pattern means.

\paragraph{Findings.} Across \nModels\ models from \nFamilies\ families and \nEpisodes\ negotiations:
\begin{itemize}
  \item \textbf{Obedient agents.} Agents almost never broke the client's requirement or leaked its confidential figure (\violationCount\ violations, \leakCount\ leaks).
  \item \textbf{Silent reports.} Reports conveyed the blocked better deal in a small minority of exposed episodes (judged non-disclosure \doneJudgedAll\ of \exposedCount\ exposed episodes). \doneLowerAll\ of reports referred to no other offer or option at all, which is a judge-free floor, and \doneNotNamed\ never gave the offer's price, points or blocked term.
  \item \textbf{Knowledge without disclosure.} Asked directly afterwards, agents whose reports had been silent named the blocked offer's value correctly in \probeKnewAll\ of cases. When the client's question asked about turned-down terms, judged non-disclosure fell from \intOpenJudged\ to \intTradeoffsJudged.
  \item \textbf{Capability and disclosure.} On the judge-free measure, non-disclosure fell with capability rank within every family that has a ladder, so H4 was not supported. On the judged measure the fall was confined to one developer: Claude Opus 5 and Claude Fable 5.1 left the blocked deal unconveyed in \doneJudgedClaudeOpusFive\ and \fableJudged\ of exposed episodes, against \doneJudgedGptSixAstra\ for GPT-6 Astra, which negotiated equally well.
  \item \textbf{Capability and accuracy.} Errors in the stated value of the deal were concentrated in less capable models. Claude Haiku 4.5 misstated its deal's value in \dtwoWrongClaudeHaikuFourFive\ of the reports that stated one, against \dtwoWrongGptSixAstra\ for GPT-6 Astra and \dtwoWrongGrokFourSix\ for Grok 4.6.
\end{itemize}

\paragraph{Contributions.}
\begin{enumerate}
  \item A testbed with exact ground truth for what a delegated agent should tell its principal, in which the principal-relevant fact appears identically for every model and condition.
  \item Two judge-free measures of report fidelity, blocked-alternative non-disclosure (\Done) and stated-value error (\Dtwo), with validated model judgments as a second stage.
  \item A study of \nModels\ models under a prespecified analysis plan, with outcome quality, audit framing and model family varied on identical negotiations, plus a knowledge probe that separates not tracking from not telling.
  \item Evidence on writer-side remedies: an audit sentence, a deployer instruction, and a client prompt that asks.
\end{enumerate}
